% Structured Pruning and Post-Training Quantization for Lightweight Vision Transformers
\documentclass[conference]{IEEEtran}

\usepackage{cite}
\usepackage{amsmath,amssymb,amsfonts}
\usepackage{graphicx}
\usepackage{textcomp}
\usepackage{xcolor}
\usepackage{booktabs}
\usepackage{tabularx}

\title{Structured Pruning and Post-Training Quantization\\for Lightweight Vision Transformers on Edge Devices}

\author{
    \IEEEauthorblockN{Fatih\IEEEauthorrefmark{1},
    Oddy Virgantara Putra\IEEEauthorrefmark{2}}
    \IEEEauthorblockA{
        \IEEEauthorrefmark{1}Informatics Department, ...
    }
}

\begin{document}
\maketitle

\begin{abstract}
Vision Transformers (ViTs) have achieved remarkable performance across
various computer vision tasks, but their computational demands remain a
barrier to edge deployment. This paper systematically investigates the
compounding effect of structured pruning and post-training quantization
(PTQ) on lightweight ViT variants. We evaluate DeiT-Tiny under various
pruning ratios (25--50\%) combined with INT8 dynamic quantization. Our
results show that combining 50\% head+MLP pruning with PTQ reduces model
size from 21.9MB to 3.4MB (84.4\% reduction) while cutting FLOPs by 54\%,
demonstrating viability for resource-constrained edge deployment.
\end{abstract}

\section{Introduction}
Vision Transformers (ViTs)~\cite{dosovitskiy2020vit} have become the dominant
architecture for computer vision, but their large parameter count and high
computational cost hinder deployment on edge devices. While
quantization~\cite{quant2022pytorch} and pruning~\cite{zhu2023pruning} are
well-studied individually, their compounding effect on tiny ViTs remains
underexplored. We address this gap.

\section{Methodology}
We apply structured pruning to DeiT-Tiny~\cite{touvron2021deit}:
\begin{itemize}
    \item \textbf{Head pruning:} Remove entire attention heads by
    truncating the QKV projection matrix and updating \texttt{num\_heads}.
    \item \textbf{MLP pruning:} Reduce the intermediate dimension of
    feed-forward networks by slicing \texttt{fc1} and \texttt{fc2} weights.
    \item \textbf{Dynamic PTQ:} Apply INT8 dynamic quantization
    (\texttt{torch.ao.quantization.quantize\_dynamic}) on CPU after pruning.
\end{itemize}
Ten configurations are benchmarked covering head ratios 1.0, 0.75, 0.5
and MLP ratios 1.0, 0.75, 0.5 with and without PTQ.

\section{Results}
Table~\ref{tab:results} summarizes the key findings.

\begin{table}[ht]
\centering
\caption{DeiT-Tiny pruning + PTQ results}
\label{tab:results}
\small
\begin{tabular}{lrrrr}
\toprule
Config & Params(M) & Size(MB) & FLOPs(M) & TP(img/s) \\
\midrule
Baseline         & 5.72 & 21.87 & 1074.85 & 194 \\
H0.75\_M0.75     & 4.24 & 20.27 & 784.36  & 279 \\
H0.50\_M0.50     & 2.76 & 18.58 & 493.87  & 178 \\
PTQ-only         & 5.72 & 6.26  & 1074.85 & 102* \\
H0.75\_M0.50\_Q  & 3.35 & 3.98  & 610.07  & 145* \\
H0.50\_M0.50\_Q  & 2.76 & 3.41  & 493.87  & 162* \\
\bottomrule
\multicolumn{5}{l}{*PTQ latency on CPU; others on CUDA. TP = Throughput} \\
\end{tabular}
\end{table}

The combination of 50\% head+MLP pruning with PTQ achieves 84.4\% model
size reduction (21.9MB $\to$ 3.4MB) and 54.1\% FLOPs reduction, making
DeiT-Tiny feasible for deployment in sub-4MB flash budgets typical of
MCU-class devices.

\section{Conclusion}
Structured pruning and post-training quantization exhibit compounding
benefits for tiny ViTs. The combined pipeline reduces model size by
5--6$\times$ while halving computational cost, all without
fine-tuning or retraining. Future work will extend this to additional
architectures and evaluate on physical edge hardware.

\bibliographystyle{ieeetr}
\bibliography{refs}
\end{document}
